\subsection{大小核调度与负载均衡}\label{sec:sched}

RK3588 不是一块由八个相同核心拼成的处理器，它把两种规格的核混在同一块芯片上，四个 Cortex-A55
小核运行在 1.8\,GHz，四个 Cortex-A76 大核运行在 2.256\,GHz。大核的微架构更宽、流水线更深，
同样一段代码在大核上跑得明显比在小核上快。这种异构布局通常称作大小核（big.LITTLE），它的好处是
让轻负载落在省电的小核而把重活交给大核，可一旦操作系统不区分这两类核，把一个计算密集的线程随手
放到一个小核上，它就只能以小核的速度慢慢跑，旁边空着的大核帮不上忙。对一条每帧都要在限定时间内
跑完采集、解码、缩放与推理的感知流水线来说，线程被放到哪个核上，直接决定了这一帧能否按时交付。

Starry 继承自 ArceOS 的任务调度在这一点上采取的是最朴素的策略。一个线程被创建或被唤醒时，调度器
默认把它放到当前正在执行放置动作的那个核上，借此让它的缓存保持温热，只有当线程的亲和性掩码不允许
留在本核时，才退而用一个简单的轮询计数去挑下一个允许的核。这条路径在 \texttt{select\_run\_queue}
里就是先看 \texttt{task.cpumask().get(current\_cpu)} 是否成立，成立就直接落在本核，整个过程不读取
任何一个核当前有多忙，也不区分这个核是大核还是小核。机器人程序起初没有设置任何 CPU 亲和性，于是
采集、解码、缩放与推理这几个本可以分散开的线程全都被放到了最先启动它们的那个 A55 小核上，在
cpu0 上首尾相接地串行执行，彼此争抢同一个最慢的核。这正是迁移到 Starry 之初端到端帧时间偏高的
根源，后续把推理线程手动钉到一个 A76 大核上得到的数倍提速，以及那一组随核位变化的实测数字，都
记在性能优化一节（\ref{sec:opt}）里。

\paragraph{让调度器区分大核与小核}
\textbf{针对这一缺口，本队做了一个对大小核有感知的负载均衡器原型，它不替换原有的调度算法，而是在放置
决策与空闲核两处补上以核能力为权重的均衡，并以一个默认关闭的 \texttt{sched-loadbalance} 特性开关
控制，开关不打开时默认构建与 QEMU 持续集成的产物逐字节不变。}它要回答的第一个问题是每个核到底有
多强。这一信息来自设备树，\texttt{cpu\_capacities} 在系统启动时遍历设备树里的各个 \texttt{cpu@*}
节点，优先读取节点上的 \texttt{capacity-dmips-mhz} 属性，读不到时再退而看节点的 \texttt{compatible}
字符串，认出 \texttt{arm,cortex-a55} 就记为 530，认出 \texttt{arm,cortex-a76} 就记为 1024，两者都
没有时给一个 1024 的默认值，结果按逻辑核号缓存下来。这样大核与小核之间约一比二的算力差距，就以
一个归一化的能力权重被调度器读到，而在所有核都相同的同构板子或 QEMU 上，这套权重退化为全相等，
均衡器随之退化为单纯的负载分摊。

\paragraph{按能力加权的负载去挑核}
有了每个核的能力，放置决策就不再盯着当前核，而是去找加权后最轻的那个核。衡量轻重的不是裸的就绪
任务数，而是把它除以核能力再乘以一个固定基准得到的有效负载，即 \texttt{load * 1024 / capacity}。
同样一份待跑的工作放在小核上要花更久，除以更小的能力之后它的有效负载就更高，于是
\texttt{select\_least\_loaded} 在遍历亲和性允许且已经上线的各个核时，会自然地先把计算往大核上堆，
等大核被填满、它们的有效负载追上小核了，工作才开始往小核溢出。挑核时对线程上一次所在的核留了一点
偏向，把它的负载当作低一档来比较，使在负载相当的情况下线程尽量回到缓存温热的那个核。线程被唤醒时
走的也是同一套加权择核，只是把上次所在的核作为这点偏向的对象。除了放置这一头，空闲的核也不再立刻
停下来等中断，而是在 \texttt{idle\_pull\_once} 里扫一遍其它在线的核，从加权负载最重的那个核上偷一个
既符合目标核亲和性、又没有正在别处运行、也不是空闲任务的就绪线程过来自己跑；定时器节拍上还有一个
限频且带迟滞的反向推送，本核明显比最轻的远端核更忙时主动推一个线程过去，避免只靠偷取在某些情形下
反应不及。这些跨核搬运都严格遵守同一时刻至多持有一把调度器锁的纪律，被搬运的线程在两个队列之间的
空档里只由一个本地引用持有，因而既不会丢失也不会被两个核同时运行。

\paragraph{原型的定位与局限}
这个均衡器在 QEMU 的多核环境里通过了一组覆盖亲和性迁移、抢占与并发竞争的测试，相关的就绪队列
计数与可窃取任务的接口也都各自带了单元测试。但必须如实说明，它目前仍是一个原型，是日后把调度
真正做成对体系结构有感知的基础，而不是已经接管生产负载的调度器。对机器人当前这套以手动钉核运行的
单线程化负载而言，工作早已被显式地固定在选定的核上，再叠一层自动均衡并不改变结果，因此机器人实际
运行时仍沿用手动钉核，把这个原型发展为完整的体系结构感知调度留作后续工作。它要解决的真问题也正在
这里，跨核唤醒的延迟仍有长尾，最差的一帧约 487\,ms，而 Linux 同条件下约 22\,ms，伴随约一成的丢帧，
把这条尾巴收拢到与放置同样可控的程度，是这块调度工作尚未走完的部分。\textbf{这一原型尚未提交至上游（待提交）。}

\repfig{scheduling.png}{大小核异构下的放置：能力权重自设备树读出，放置与空闲偷取都以 \texttt{load*1024/capacity} 的有效负载择核，计算先堆向 A76 大核，填满后才向 A55 小核溢出}{fig:sched}
